%% 自動生成（phys/gen_phys.py）
\phtocgrid{\phgi{保存則と相似の地図}{4}\phgi{この本の使い方}{6}}
{\color{bkLine}\hrule height 0.35pt}\vspace{2.4mm}
\phtocchap{波-0}{ゆれは止められる？　波は使える？}{地震のゆれ、楽器の音、光。どこが同じ？}{}{3}{3}{8}{J}
\phtocpart{第1部}{ゆれる：振動と制御}{12}{1}
\phtocchap{波-1}{ばねとおもりで考える：1自由度の振動}{ブランコの1往復の時間は、何で決まる？}{6}{4}{5}{14}{C}
\phtocchap{波-2}{ゆすり方で変わるゆれ：強制振動と共振}{洗濯機は、なぜ回り始めにガタガタゆれることがある？}{3,6,9}{5}{5}{32}{C}
\phtocchap{波-3}{つながるゆれ：多自由度振動とラグランジュの式}{高いビルの屋上の重りは、どうやってゆれをおさえる？}{6,7,8}{5}{5}{50}{C}
\phtocchap{波-4}{ねらいどおりに動かす：フィードバック制御}{エアコンは、どうやって部屋をちょうどよい温度に保つ？}{9}{5}{5}{68}{C}
\phtocpart{第2部}{伝わる：波}{86}{2}
\phtocchap{波-5}{波を式で書く：進行波と波動方程式}{スタジアムのウェーブでは、何が進んでいる？}{6,8}{4}{5}{88}{C}
\phtocchap{波-6}{はね返って重なる：反射・定在波・固有振動}{ギターの弦は、なぜおさえる場所で音の高さが変わる？}{6,9}{5}{5}{106}{C}
\phtocbreak
\phtocpartcont{第2部}{伝わる：波}{2}
\phtocchap{波-7}{波が運ぶエネルギー：インピーダンスと反射}{超音波の検査で、なぜおなかにゼリーをぬる？}{6}{5}{5}{124}{C}
\phtocchap{波-8}{重なる波：干渉・回折・ドップラー効果}{救急車のサイレンは、なぜ通りすぎると低く聞こえる？}{4}{4}{4}{142}{C}
\phtocpart{第3部}{はかる：音・振動の計測と信号}{158}{3}
\phtocchap{波-9}{音の大きさをはかる：音速・デシベル・騒音}{雷が光ってから、なぜ少しおくれて音が聞こえる？}{3}{4}{4}{160}{C}
\phtocchap{波-10}{見えない中を調べる：超音波と振動の計測}{コウモリは、暗やみでどうやってものの位置を知る？}{9}{4}{4}{176}{C}
\phtocchap{波-11}{音を数字にする：信号とスペクトル}{動画の車輪が、逆に回って見えることがあるのはなぜ？}{3,9}{5}{5}{192}{C}
\phtocpart{第4部}{曲がる：光}{210}{4}
\phtocchap{波-12}{光の進み方：反射・屈折・レンズ}{水に入れたストローは、なぜ折れて見える？}{5,7}{4}{4}{212}{C}
\phtocchap{波-13}{光の波：干渉・回折・偏光・散乱}{空は、なぜ青い？　夕焼けは、なぜ赤い？}{4,9}{4}{5}{228}{C}
\phtocchap{波-14}{波で情報を送る}{スマホの声は、どうやって電波にのせて送る？}{9}{4}{5}{246}{R}
\par\vspace{1.5mm}
\phtocgrid{\phgi{こたえ}{254}\phgi{記号カード}{261}\phgi{物性値・定数・単位}{264}\phgi{\mbox{なぜ？}の道すじ}{268}\phgi{できるようになったことチェック表}{270}\phgi{つぎの巻へ}{272}\phgi{さくいん}{274}}
